% songs/die-ballade-vom-pfeiffer.tex
\begin{song}{Die Ballade vom Pfeiffer}{Reinhard Mey}{Capo -}

\cl{      Am                      E7            Am}
\slbl{1.}Ich ritt aus San Alfredo im letzten Tageslicht,
\cl{    C               C5                             Am}
Vielleicht auch aus El Paso, so genau weiß ich das nicht.
\cl{                                    D              D4      D}
Seit vierzig Tagen, vierzig Nächten war ich auf der Flucht,
\cl{   Am                                  E7           Am}
In jeder Stadt fand ich mein Bild, darunter stand: GESUCHT!
\cl{    G                C                G             C}
Sie nannten mich den Stillen, und man flüsterte mir nach,
\cl{     G               C                 D             E7}
Daß, wenn ich was zu sagen hätte, mein Colt für mich sprach.

\slbl{Refrain: (pfeifen)}
\cl{Am   Em   Am   G   Am}

\cl{      Am                        E7               Am}
\slbl{2.}Sie nannten mich den Pfeifer, und meine Devise hieß:
\cl{     C               C5                                  Am}
Wenn Dir wer vor die Mündung kommt, erst pfeife und dann schieß!
\cl{                                          D             D4       D}
Warum ich pfiff, das weiß ich nicht, weiß nicht mal wie ich heiß,
\cl{   Am                              E7           Am}
Im Westen ist es niemals gut, wenn einer zuviel weiß!
\cl{    G                 C                G              C}
Ich weiß nur, wo mein Lied erklang, da wurden Bretter knapp,
\cl{     G               C             D          E7}
Weil jeder Schreiner wußte, daß es Arbeit für ihn gab.

\cl{      Am                          E7                 Am}
\slbl{3.}Als ich nach Bloody Corner kam, sah ich von weitem her:
\cl{    C           C5                              Am}
Die Summe unter meinem Namen hatte zwei Stellen mehr.
\cl{                                      D           D4        D}
Ein Prämienjäger sagte: "Pfeifer, ich wart' schon auf dich!"
\cl{    Am                                      E7        Am}
Ich fuhr herum, pfiff einen Ton dann sprach mein Colt für mich.
\cl{    G               C           G               C}
Als wenig später im Saloon mein Achtundreiß'ger spie,
\cl{   G           C               D        E7}
Da spielte der Mann am Klavier dazu die Melodie.

\cl{      Am                      E7               Am}
\slbl{4.}Die Ellenbogen aufgestützt, die Flügeltür im Blick
\cl{      C             C5                             Am}
Stand neben mir ein Fremder, sehr glattrasiert und dick.
\cl{                                  D             D4       D}
"Man nennt mich hier den Denker", so stellte er sich vor,
\cl{     Am                          E7        Am}
spie lässig in den Spucknapf und sagte mir ins Ohr:
\cl{      G              C                     G              C}
"Wenn Du zehntausend Dollar brauchst, dann hab' ich einen Plan,
\cl{G                 C         D            E7}
Todsicher, genial einfach!" Dankbar nahm ich an.

\cl{        Am                        E7           Am}
\slbl{5.}Jetzt sitz' ich hinter Gittern, von Zweifeln angenagt,
\cl{C                       C5                             Am}
Vielleicht war doch des Denkers Plan so gut nicht, wie er sagt.
\cl{                                      D            D4        D}
Er sagte: "Das bringt Dir zehntausend Dollar, wenn Du's wagst
\cl{    Am                              E7                 Am}
Zum Sheriff ins Büro zu gehen, Dich vorstellst und ihm sagst,
\cl{     G                 C                  G               C}
Grüß Gott, ich bin der Pfeifer, ich komm' selber, wie ihr seht,
\cl{   G                C                  D           E7}
Um die Belohnung zu kassier'n, die auf meinem Kopf steht!"

\cl{      Am                              E7        Am}
\slbl{6.}Ich sitz' auf meines Pferdes Rücken unter dem Galgenbaum,
\cl{      C                C5                               Am}
Einen Strick um meinen Hals, der Henker hält mein Pferd im Zaum.
\cl{                                      D             D4     D}
Gleich gibt er ihm die Zügel und dann ist's mit mir vorbei,
\cl{    Am                              E7               Am}
Der Totengräber gräbt mein Grab und pfeift mein Lied dabei.
\cl{    G                 C                  G                C}
Der Scharfrichter tut seine Pflicht mein Pferd setzt sich in Trab
\cl{    G                C}
Und unten brüllt der Regisseur:
"Verdammte Schlamperei, jetzt ist uns der Ast zum dritten Mal abgebrochen!
Der Film ist auch gerissen. Also Kinder, für heute ist Feierabend,
die Leiche dreh'n wir morgen ab!" \cl{[D E7 Am]}
\griffe{C5, D4}
\end{song}
